Let $x_1 \sim \Poiss(\lambda_1)$ and $x_2 \sim \Poiss(\lambda_2)$ be two
independent Poisson random variables, that is,
$p(x_i = k) = \exp(-\lambda_i) \frac{\lambda_i^k}{k!}$,
$k = 0, 1, 2, \ldots$ Prove that
$x_1 + x_2 \sim \Poiss(\lambda_1 + \lambda_2)$.

\begin{solution}
\solpart{Idea}
To find the distribution of a function $y = f(x_1, \ldots, x_n)$ of discrete
random variables, find $p(y = k)$ for each value $k$:
\begin{enumerate}
  \item Write the event $\{f(x_1, \ldots, x_n) = k\}$ as a union of
    disjoint events $\{x_1 = j_1, \ldots, x_n = j_n\}$, one for each tuple
    with $f(j_1, \ldots, j_n) = k$.
  \item The probability of a union of disjoint events is the sum of their
    probabilities:
    \[
      p(y = k) = \sum_{f(j_1, \ldots, j_n) = k}
        p(x_1 = j_1, \ldots, x_n = j_n).
    \]
  \item If the variables are independent, write each joint probability as
    a product.
  \item Simplify the sum and compare it with known distributions.
\end{enumerate}

\solpart{Solution}
\solstep{Step 1. Write the event as a union of disjoint events.}
The variables take only the values $0, 1, 2, \ldots$ Thus $x_1 + x_2 = k$ if
and only if $x_1 = j$ and $x_2 = k - j$ for one $j \in \{0, \ldots, k\}$:
\[
  \{x_1 + x_2 = k\} = \bigcup_{j=0}^{k} \{x_1 = j, \ x_2 = k - j\}.
\]
The events are disjoint, because each event has a different value of $x_1$.

\solstep{Step 2. Sum the probabilities of the disjoint events.}
\[
  p(x_1 + x_2 = k) = \sum_{j=0}^{k} p(x_1 = j, \ x_2 = k - j).
\]

\solstep{Step 3. Use the independence of $x_1$ and $x_2$.}
\[
  p(x_1 + x_2 = k) = \sum_{j=0}^{k} p(x_1 = j) \, p(x_2 = k - j).
\]

\solstep{Step 4. Substitute the Poisson probabilities.}
\[
  p(x_1 + x_2 = k)
  = \sum_{j=0}^{k} e^{-\lambda_1} \frac{\lambda_1^j}{j!} \cdot
    e^{-\lambda_2} \frac{\lambda_2^{k-j}}{(k-j)!}.
\]

\solstep{Step 5. Move the factors that do not depend on $j$ out of the sum.}
\[
  p(x_1 + x_2 = k)
  = e^{-\lambda_1} e^{-\lambda_2}
    \sum_{j=0}^{k} \frac{\lambda_1^j \lambda_2^{k-j}}{j! \, (k-j)!}
  = e^{-(\lambda_1 + \lambda_2)}
    \sum_{j=0}^{k} \frac{\lambda_1^j \lambda_2^{k-j}}{j! \, (k-j)!}.
\]

\solstep{Step 6. Multiply and divide by $k!$.}
\[
  p(x_1 + x_2 = k)
  = \frac{e^{-(\lambda_1 + \lambda_2)}}{k!}
    \sum_{j=0}^{k} \frac{k!}{j! \, (k-j)!} \lambda_1^j \lambda_2^{k-j}
  = \frac{e^{-(\lambda_1 + \lambda_2)}}{k!}
    \sum_{j=0}^{k} \binom{k}{j} \lambda_1^j \lambda_2^{k-j}.
\]

\solstep{Step 7. Use the binomial theorem.}
The binomial theorem gives
$\sum_{j=0}^{k} \binom{k}{j} \lambda_1^j \lambda_2^{k-j}
= (\lambda_1 + \lambda_2)^k$. Thus
\[
  p(x_1 + x_2 = k)
  = e^{-(\lambda_1 + \lambda_2)} \frac{(\lambda_1 + \lambda_2)^k}{k!}.
\]
This is the probability mass function of
$\Poiss(\lambda_1 + \lambda_2)$.

\solpart{Conclusion}
The sum of two independent Poisson random variables is a Poisson random
variable, and the parameters add. By induction, the same is true for any
finite number of independent Poisson random variables.
\end{solution}
